\chapter{Introduzione allo Sviluppo Mobile e Android}

\section{Il mercato Mobile e lo Sviluppo Cross-Platform}
Oggi il mercato dei sistemi operativi mobili è dominato da due attori principali: \textbf{Android} (Google) che detiene circa il 70-74\% del mercato globale, e \textbf{iOS} (Apple) con il restante 25-30\%.

Quando si decide di sviluppare un'applicazione mobile, ci si trova davanti a un bivio architetturale:
\begin{itemize}
    \item \textbf{Sviluppo Nativo:} Scrivere codice specifico per ogni piattaforma (Kotlin/Java per Android, Swift/Objective-C per iOS).
    Garantisce le massime performance e accesso immediato alle API del dispositivo.
    \item \textbf{Sviluppo Cross-Platform:} Scrivere il codice una sola volta e deployarlo su più piattaforme.
    Abbassa i costi e i tempi di sviluppo, ma può portare a compromessi sulle prestazioni e ritardi nel supporto alle nuove API\@.
\end{itemize}

I framework Cross-Platform più popolari ad oggi includono:
\begin{itemize}
    \item \textbf{React Native (Facebook):} Basato su JavaScript.
    \item \textbf{Flutter (Google):} Basato sul linguaggio Dart, rendering custom (non usa le UI native ma le ridisegna).
    \item \textbf{Kotlin Multiplatform Mobile (KMM):} Creato da JetBrains.
    Non è un framework UI, ma un \textit{SDK} che permette di condividere la logica di business (networking, database, logica)
    in Kotlin, mantenendo lo sviluppo della UI strettamente nativo per Android e iOS.
\end{itemize}

\section{Architettura di un'App Android e Componenti Core}
Un'applicazione Android non ha un singolo ``entry point'' (come il classico \texttt{main()} in C o Java desktop), ma è composta da quattro blocchi fondamentali (Componenti):
\begin{enumerate}
    \item \textbf{Activities:} Rappresentano le singole schermate con un'interfaccia utente (UI).
    \item \textbf{Services:} Eseguono operazioni di lunga durata in background, senza UI (es.
    riproduzione musicale, download).
    \item \textbf{Broadcast Receivers:} Ascoltano e reagiscono a messaggi o eventi di sistema (es.
    "batteria scarica", "schermo sbloccato").
    \item \textbf{Content Providers:} Gestiscono e condividono un set di dati tra diverse applicazioni in modo sicuro (es.
    la rubrica dei contatti).
\end{enumerate}

Lo sviluppo moderno (Modern Android Development) incoraggia l'uso di librerie \textbf{Jetpack} (es. ViewModel per lo stato, WorkManager per il background, Room per i DB) che semplificano enormemente la gestione del ciclo di vita di questi componenti.

\section{Android Studio e il processo di Build (Gradle)}
\textbf{Android Studio} è l'IDE ufficiale di Google basato su IntelliJ IDEA. Dal 2015 ha sostituito Eclipse.

\subsection{Gradle e Architettura Hardware}
Nei quiz è importante ricordare che:
\begin{itemize}
    \item \textbf{Gradle} è un sistema di build automation open-source.
    Costruito sui concetti di Apache Ant e Maven, utilizza un \textbf{DSL (Domain-Specific Language)} basato su Groovy o Kotlin.
    \item I dispositivi mobili utilizzano spesso l'\textbf{architettura ARM big.LITTLE}: un'architettura che combina
    nuclei di processori di dimensioni diverse per \textbf{ottimizzare le prestazioni e l'efficienza energetica}.
\end{itemize}

\subsection{Dall'Idea all'Esecuzione: APK, DEX e ART}
Nei quiz è fondamentale conoscere come il codice diventa un'app eseguibile:
\begin{enumerate}
    \item Il codice sorgente (Kotlin/Java) viene compilato nel bytecode standard di Java (\texttt{.class}).
    \item Il compilatore Android converte i file \texttt{.class} nel formato \textbf{DEX (Dalvik Executable)}, ottimizzato per dispositivi mobili.
    \item Il tool \texttt{aapt} unisce i file DEX e le risorse compilate per creare l'archivio \textbf{APK} (Android Application Package).
\end{enumerate}

\begin{tcolorbox}[colback=blue!5!white,colframe=blue!75!black,title=Android Runtime (ART) vs Dalvik]
    In passato, Android usava la macchina virtuale \textbf{Dalvik} (compilazione JIT - Just In Time).
    Oggi utilizza \textbf{ART (Android Runtime)}.
    ART compila l'app in codice macchina nativo al momento dell'installazione tramite la tecnica \textbf{AOT (Ahead-Of-Time) compilation}.
    \newline\textit{Vantaggi:} Esecuzione molto più veloce. \newline\textit{Svantaggi:} Tempi di installazione leggermente più lunghi e maggiore spazio occupato sul disco.
\end{tcolorbox}

\section{Figma: UI vs UX e Prototipazione}
Prima di scrivere il codice, la fase di design è fondamentale.
\begin{itemize}
    \item \textbf{UX (User Experience):} Struttura logica, gerarchie, posizionamento e flussi funzionali.
    Analizza i comportamenti reali degli utenti.
    \item \textbf{UI (User Interface):} Struttura visiva, colori, font, icone.
    Assicura che il prodotto sia visivamente stimolante.
\end{itemize}

\begin{tcolorbox}[colback=yellow!10!white,colframe=yellow!70!black,title=Figma nei Quiz]
    \textbf{Figma} è una piattaforma di design collaborativo \textbf{basata su cloud} (Web-based tramite WebGL/WASM).
    Utilizza i \textbf{CRDT} (Conflict-free Replicated Data Types) e \textbf{WebSockets} per permettere a più persone
    di lavorare sullo stesso documento in \textit{real-time} senza conflitti di file (approccio \textit{Last-Writer-Wins}).
    Consente di creare \textbf{veri prototipi} per testare l'esperienza utente prima di scrivere codice.
    \begin{itemize}
        \item \textbf{Caratteristica principale:} È \textit{Web-based} (WASM + WebGL).
        A differenza di software storici come Sketch, non richiede installazioni locali o file \texttt{.sketch} sul disco.
        \item Il motore di rendering sotto il cofano è scritto in \textbf{C++} e transpilato in WebAssembly, permettendo la
        gestione di grafica vettoriale con la potenza di un videogioco direttamente nel browser.
    \end{itemize}
\end{tcolorbox}

\subsection{Il contenuto di un APK}
Da un punto di vista dei file, un APK è un archivio ZIP che contiene:
\begin{itemize}
    \item \texttt{classes.dex}: Il codice eseguibile dell'applicazione.
    \item \texttt{res/}: Le risorse non compilate.
    \item \texttt{resources.arsc}: Le risorse compilate (tabelle di lookup generate da \texttt{aapt}).
    \item \texttt{META-INF/}: Informazioni sulle firme digitali.
    \item \textbf{\texttt{AndroidManifest.xml}}: Il file fondamentale che descrive l'app al sistema operativo Android.
\end{itemize}

\section{AndroidManifest.xml}
Ogni app Android deve avere un file chiamato \texttt{AndroidManifest.xml} nella root del progetto. 
Definisce le caratteristiche essenziali dell'applicazione necessarie ad Android prima ancora che l'app venga avviata:
\begin{itemize}
    \item Dichiara il \textit{package name} (che funge da ID univoco dell'app).
    \item \textbf{Dichiara i Componenti:} Tutte le \texttt{<activity>}, \texttt{<service>}, \texttt{<receiver>} e \texttt{<provider>} devono essere dichiarati qui. Se crei una classe Activity in Kotlin ma non la dichiari nel Manifest, l'app andrà in crash se cerchi di avviarla.
    \item Dichiara i \textbf{permessi} (\texttt{<uses-permission>}) richiesti (es. INTERNET, fotocamera).
    \item Dichiara le caratteristiche hardware richieste (es. touch screen, NFC).
\end{itemize}

Un costrutto chiave all'interno del manifest è l'\textbf{Intent-Filter}. Per indicare al sistema operativo qual è la prima Activity da lanciare all'avvio dell'app (l'Activity \textit{Home}), si usa la seguente sintassi:
\begin{lstlisting}[language=XML]
<activity android:name=".MainActivity" android:exported="true">
    <intent-filter>
        <action android:name="android.intent.action.MAIN" />
        <category android:name="android.intent.category.LAUNCHER" />
    </intent-filter>
</activity>
\end{lstlisting}
Il flag \texttt{android:exported="true"} indica che questa Activity può essere avviata da app esterne (incluso il launcher (la schermata home) del sistema operativo).